\section{Данные и протокол оценки}

\paragraph{Датасет.} Yambda~\cite{yambda} в конфигурации 50M: 9\,238
пользователей, 872\,861 айтем с хотя бы одним прослушиванием, 46.0 млн
прослушиваний и 872 тыс.\ лайков в обучающем периоде. События содержат метку
органического прослушивания, долю прослушивания \code{played\_ratio\_pct}
и длину трека. Из всех прослушиваний 64.7\% дослушаны не менее чем наполовину;
распределение популярности айтемов имеет коэффициент Джини 0.907, то есть
каталог сильно скошен в сторону верхушки.

\paragraph{Протокол Listen+.} Как и в официальном бенчмарке
(\code{TRACK\_LISTEN\_THRESHOLD = 50}), прослушивание попадает в историю,
словарь и таргеты только при \code{played\_ratio\_pct} $\ge 50$. В обучающем
периоде это 29\,135\,186 событий (63.3\% всех прослушиваний) у 9\,207
пользователей; средняя доля прослушивания в них 97.3\%. Медианная история
пользователя --- 1\,777 событий, 90-й перцентиль --- 8\,028. Отсечённые
прослушивания (доля досмотра ниже половины) модель никогда не видит, поэтому
сигнал «трек отклонён» доступен только reward model как отдельная ветка
кандидатов (раздел~\ref{sec:method-rm}).

\paragraph{Разбиение.} Используется разбиение поставщика данных, совпадающее
с глобальным временным сплитом официального бенчмарка: обучение до
$t=25\,823\,595$ с, 30-минутный разрыв, день валидации
($25\,825\,400 \le t < 25\,911\,795$) и день теста
($25\,913\,600 \le t < 26\,000\,000$). Границы ETL вычисляет из констант
\code{last\_timestamp = 26\,000\,000}, \code{test\_size = 86\,400},
\code{val\_size = 86\,400}, \code{gap\_size = 1\,800} и проверяет
фактические минимумы и максимумы сплитов: условие
«train max ts $<$ test min ts» выполняется и пишется в \code{eda.json}.

\paragraph{Словарь.} Словарь Listen+ строится по всем трём сплитам и содержит
630\,974 айтема (72.3\% каталога всех прослушиваний): 627\,648 встречаются
в обучении, 3\,326 --- только в валидации и тесте. Айтемы вне обучающего
Listen+ остаются в словаре и достижимы при ранжировании, как и в официальном
коде (\code{drop\_non\_train\_items = False}); их эмбеддинги инициализируются
из общего распределения, но не обучаются. Из тестовых целей 1\,782 айтема
(3.24\%) не встречаются в обучающем Listen+; это 2\,332 события (1.48\%)
у 688 пользователей (14.96\%). Отсюда практический потолок Recall@100 ---
98.52\%.

\paragraph{Тестовый день.} 232\,049 событий, 4\,812 пользователей; после
фильтра Listen+ --- 157\,696 целей у 4\,600 пользователей, в среднем 34
цели на пользователя. Лайки тестового дня: 4\,279 событий у 1\,311
пользователей, 97.3\% из них новые (не повторный лайк уже лайкнутого
айтема). В like-метрику входят 1\,098 пользователей (3\,445 лайков):
200 пользователей (15.3\%) не имеют истории Listen+ и выпадают из оценки,
115 айтемов из 3\,739 лежат вне словаря.

\paragraph{Метрики.} NDCG@$k$, Recall@$k$ и MAP@$k$ для
$k\in\{10,50,100\}$, Coverage@$k$ --- доля каталога, попавшая в выдачи.
NDCG@$k$ нормируется на идеальный DCG с $\min(|\text{цели}|, k)$
релевантными позициями, MAP@$k$ --- на то же число. Оценка идёт по полному
каталогу: модель считает логиты для всех 630\,974 айтемов порциями по
100\,000 и отдаёт top-100, поэтому метрики не смещены выбором кандидатов.
Отдельно считаются like-NDCG@10 и like-Recall@100 по лайкам тестового дня;
это единственная метрика, которая не является выходом reward model,
оптимизируемой RL. Официальная Like-задача --- likes-to-likes (история из
последовательности лайков), поэтому её опубликованные значения служат
ориентиром масштаба, а не прямым baseline.

\paragraph{Единица анализа.} Пользователь, а не показ и не событие. Все
сравнения базовой и RL-политики --- парные по пользователям, доверительные
интервалы считаются бутстрэпом (10\,000 ресэмплов).

\paragraph{Ориентиры.} Помимо популярностного baseline, в таблицах приведены
опубликованные значения официального бенчмарка для SASRec~\cite{sasrec} на
Listen+: NDCG@10 0.0748, Recall@100 0.1026. Мы их не переобучали и
сравниваем с нашим прогоном как с внешней точкой отсчёта.
